\documentclass[conference]{IEEEtran}
\usepackage{cite}
\usepackage{amsmath,amssymb,amsfonts}
\usepackage{booktabs}
\usepackage{array}
\usepackage{url}
\usepackage{xcolor}
\usepackage{graphicx}
\usepackage{multirow}

\newcommand{\designonly}{\textsc{Design-Only}}
\newcommand{\wpp}{\textit{World Population Prospects 2024}}
\newcommand{\lpd}{\mathrm{LPD}}
\newcolumntype{L}[1]{>{\raggedright\arraybackslash}p{#1}}
\setlength{\emergencystretch}{1em}

\begin{document}

\title{Will the World's Population Keep Growing?\\
A Component-Resolved, Calibrated Proposal for Population-Peak Persistence}

\author{\IEEEauthorblockN{Anonymous Research Plan Author}
\IEEEauthorblockA{\textit{Proposal generated from a frozen Survey--Idea--ExperimentDesign record}\\
Design-only scientific research plan}}

\maketitle

\begin{abstract}
Whether world population will continue to grow is often expressed as a request for one number or one peak year. That framing discards the demographic processes and structural uncertainty that create projections. This paper presents a proposal, not an executed forecast, for a component-resolved Population Peak Persistence Atlas. It will propagate uncertainty in country-by-age-by-sex fertility, survival, and net migration through a cohort-component accounting system, aggregate posterior trajectories to the world level, and estimate the conditional probability that annual global population growth remains positive at 2100. A demographic-components-only baseline will be compared with an otherwise matched model containing pre-registered climate-disruption features. Rolling-origin validation, interval calibration, predictive scoring, component-swap ablations, and aggregation checks define the acceptance criteria. The proposal neither asserts a single future population total nor treats climate variables as a universal causal determinant. Its contribution is a falsifiable protocol that separates what is conditionally supported by a calibrated model from what remains structurally uncertain and requires human review.
\end{abstract}

\begin{IEEEkeywords}
population projection, demographic forecasting, cohort-component model, fertility, mortality, migration, calibration, uncertainty, climate risk
\end{IEEEkeywords}

\section{Introduction}

Global population change is the aggregate result of births, deaths, and international redistribution across populations with different age and sex structures. It is therefore tempting, but scientifically misleading, to reduce the question "Will the world's population keep growing indefinitely?" to a single forecasted total. A numerical trajectory may be useful for planning, yet it cannot by itself disclose which assumptions generated it, how it is calibrated, or whether a different but defensible model family yields another peak date. The question is also temporally ill posed: no finite data record can prove an infinite-horizon proposition.

The most useful operational question is narrower. Under explicitly specified trajectories of fertility, mortality, migration, and possible climate-related disruption, does a calibrated probabilistic demographic model assign positive global growth at the terminal year 2100? The corresponding scientific endpoint is a conditional probability and a distribution of first peak years. It is not a declaration that the population will inevitably peak, decline, or grow forever.

This distinction is important because credible projections can legitimately disagree. The United Nations \wpp{} dataset is the official global baseline, providing estimates from the 1950s and projections to 2100 for 237 countries or areas under a range of plausible outcomes \cite{unwpp2024}. In another verified study, Vollset \textit{et al.} modelled fertility, mortality, migration, and population for 195 countries and territories; their reference scenario peaked in 2064 at 9.73 billion and declined to 8.79 billion in 2100, with uncertainty intervals \cite{vollset2020}. These results describe that model's inputs and structure. They do not establish a universal end-of-century truth.

This research proposal develops a transparent way to compare such conditional branches. It follows a frozen four-stage provenance chain: Survey evidence and accepted gaps, an evidence-bound idea portfolio, a \designonly{} experiment design, and this Author-stage report. No forecasts, model scores, posterior draws, or empirical causal effects are reported. The paper instead specifies what would count as evidence, what would falsify the proposal's central hypothesis, and what must remain under expert review.

\section{Background, Survey, and Research Gap}

\subsection{Population change is component-based}

Population forecasts commonly use cohort-component accounting. Fertility determines the creation of new cohorts, survival determines the retention and aging of existing cohorts, and net migration redistributes people across national populations. Age structure mediates all three. A sustained period of subreplacement fertility can coexist with continuing total growth if large younger cohorts enter childbearing ages, while improved survival can enlarge older populations without generating additional births. Migration sums to zero at the fully global level, but can still affect national age structures, exposure to demographic regimes, measurement quality, and the pace at which country trajectories aggregate to a world total.

The Survey therefore treats fertility, mortality, and migration as separable but coupled uncertainty-bearing processes. This choice is consistent with the verified long-horizon forecasting study of Vollset \textit{et al.} \cite{vollset2020}, the Bayesian projection approach described by Raftery, Alkema, and Gerland \cite{raftery2014}, and the explicit treatment of migration uncertainty by Azose, Sevcikova, and Raftery \cite{azose2016}. The purpose is not to reproduce any one source's model; it is to adopt their methodological lesson that a total-population curve must be traced back to uncertainty in its components.

\subsection{Why projection disagreement is evidence, not a defect}

Gerland \textit{et al.} found, using the data and Bayesian methodology available at the time, that population stabilization was unlikely within that century \cite{gerland2014}. The later WPP 2024 framework offers updated estimates and a new collection of trajectories \cite{unwpp2024}; the 2020 forecasting study reports a still different reference peak \cite{vollset2020}. These contrasts do not allow a reader to select a favored number by authority. They show that horizon, data vintage, model structure, fertility assumptions, migration treatment, and uncertainty representation are part of the result.

Probabilistic population forecasting is designed to expose this dependence. Raftery and Sevcikova summarize its use from short to very long horizons and distinguish uncertainty-aware projections from deterministic extrapolations \cite{raftery2023}. Yet an uncertainty interval alone is not enough. A future interval can look precise while failing calibration on periods that can be checked. The proposed work consequently makes rolling-origin prediction and interval calibration gates, rather than a visually persuasive 2100 fan chart, central scientific criteria.

\subsection{Climate variables require a narrow role}

Climate disruption can plausibly interact with mortality hazards, displacement, health-care access, food systems, livelihoods, and reproductive conditions. The mechanisms are heterogeneous in time and place and are entangled with development, adaptation, conflict, measurement, and policy. The frozen Survey record does not provide a universal causal coefficient connecting a single climate index to world population. This proposal therefore refuses that shortcut.

Instead, pre-registered climate-disruption features, if approved after provenance review, will be treated as candidate lagged predictors inside demographic component submodels. The baseline is a demographic-components-only model. A climate-feature model is retained only if it improves out-of-sample predictive score and calibration under identical data cutoffs and posterior-draw budgets. This is a predictive ablation, not an estimate of a global causal effect.

\subsection{Accepted gaps}

The Survey accepted two bounded gaps. First, a transparent, component-resolved protocol is needed to compare post-2100 persistence conditions while preserving uncertainty in fertility, mortality, migration, and possible climate-disruption modifiers. Second, projection comparison needs a decision layer that decomposes peak persistence into component assumptions, calibration, and thresholds rather than presenting only competing headline totals. The Survey rejected the claim that existing methods ignore uncertainty, because verified Bayesian methods directly contradict it. It retained climate identifiability as an explicit uncertainty, not a gap to be filled by assertion.

\begin{table*}[t]
\caption{Frozen evidence anchors and their permitted use in this proposal}
\label{tab:evidence}
\centering
\footnotesize
\begin{tabular}{p{0.19\textwidth}p{0.19\textwidth}p{0.48\textwidth}}
\toprule
Source & Verified scope & Permitted use in this proposal \\
\midrule
UN WPP 2024 \cite{unwpp2024} & Official estimates and projections for 237 countries or areas through 2100 & Core accounting-data candidate and institutional projection baseline; not a claim that one central path is certain. \\
Vollset et al. \cite{vollset2020} & Component forecasting, alternative scenarios, and uncertainty intervals & Motivation for separate fertility, mortality, and migration components; its peak is reported only as that study's result. \\
Gerland et al. \cite{gerland2014} & Bayesian conclusion under an earlier data and modeling context & Evidence that conclusions can change with information and assumptions; not a contradiction to be resolved by a vote. \\
Raftery et al. \cite{raftery2014}; Azose et al. \cite{azose2016}; Raftery and Sevcikova \cite{raftery2023} & Probabilistic component and migration methodology & Design basis for uncertainty propagation, forecast scoring, and migration treatment. \\
\bottomrule
\end{tabular}
\end{table*}

\section{Research Questions and Planned Contributions}

The primary research question is: \emph{under which jointly stated fertility, mortality, migration, and optional climate-disruption trajectories does a calibrated cohort-component ensemble assign positive global population growth through 2100?} Four supporting questions make this operational.

\begin{enumerate}
\item How much of the first-peak-year distribution is attributable, conditionally, to fertility, survival, migration, and starting age structure?
\item Can a component ensemble remain calibrated at country, regional, and world aggregation levels under rolling historical holdouts?
\item Do pre-registered climate-disruption features improve prediction and calibration over a baseline built only from demographic components?
\item When do posterior persistence decisions remain stable under reasonable changes in input vintage, priors, country harmonization, and forecast origin?
\end{enumerate}

The planned contributions are methodological. First, the proposal fixes coherent endpoints: positive global growth at 2100, first global peak year, country and regional calibration, and component-swap sensitivity. Second, it uses a common cohort accounting update across all model variants so a complex model cannot win by changing the population arithmetic. Third, it makes climate information an ablation with a failure mode rather than an unexamined story. Fourth, it specifies a communication boundary: end-of-century statements remain conditional on a validated specification and cannot be converted into prescriptions about reproduction or migration.

\section{Idea Source Checkpoints and Direction Selection Audit}

The Idea stage began only after the Survey manifest was frozen. Its first candidate was a projection-reconciliation dashboard. It would make published peak dates and totals easier to compare but could not identify whether a difference arose from fertility, survival, migration, age structure, or calibration. It was not selected because it does not test either accepted gap.

The second candidate, selected as the primary direction, is the Component-Resolved Population Peak Persistence Atlas. It propagates posterior component draws through country-age-sex cohorts, aggregates them coherently, and reports conditional persistence probabilities with model diagnostics. This route directly addresses the two accepted gaps: it converts headline disagreements into testable component differences and adds a calibration-and-decision layer.

A third candidate proposed a single climate-to-population regression. It was rejected because the Survey did not verify a homogeneous, globally identifiable climate mechanism. A high-risk integrated assessment route was also retained in the portfolio but not selected: it would require strong feedback claims between demographic change, economic development, emissions, and climate outside the frozen evidence registry.

No actual Monte Carlo tree search, evolutionary optimization, or autonomous idea evolution was run. The Idea record preserves a route comparison rather than fabricating a trajectory. The selection is auditable through the target gaps, the source keys passed from Survey, the stated falsification criteria, and the feasibility path. This matters because a proposal is weakened, not strengthened, when it invents a search history that did not occur.

\section{Problem Definition, Assumptions, and Hypotheses}

\subsection{Population accounting}

Let $P_{c,a,s,t}$ be the population for country $c$, age $a$, sex $s$, and year $t$. For a non-newborn cohort, the planned accounting update is
\begin{equation}
P_{c,a+1,s,t+1}=P_{c,a,s,t}S_{c,a,s,t}+M_{c,a,s,t},
\label{eq:cohort}
\end{equation}
where $S_{c,a,s,t}$ is the annual survival probability and $M_{c,a,s,t}$ is net migration assigned to that cohort. Equation \eqref{eq:cohort} is population accounting, not an estimated effect.

The youngest cohort is generated from births,
\begin{equation}
B_{c,t}=\sum_{a\in\mathcal{A}_F}\sum_s P_{c,a,s,t} f_{c,a,s,t},
\label{eq:births}
\end{equation}
where $\mathcal{A}_F$ is the registered fertility-age set and $f_{c,a,s,t}$ is the age-specific fertility rate. Sex allocation of births and infant survival will follow a documented convention shared across all model variants.

The global total and annual growth are
\begin{equation}
G_t=\sum_c\sum_a\sum_s P_{c,a,s,t}, \qquad \Delta G_t=G_t-G_{t-1}.
\label{eq:global}
\end{equation}
Here $G_t$ is total world population in the harmonized geography and $\Delta G_t$ is annual absolute global change.

\subsection{Endpoint and assumptions}

For each posterior draw $r$, the first peak within the horizon is
\begin{equation}
\tau^{(r)}=\min\{t:G_t^{(r)}=\max_{u\leq2100}G_u^{(r)}\},
\label{eq:peak}
\end{equation}
where a draw with increasing population through 2100 is flagged as right-censored rather than assigned an invented later peak. The primary persistence endpoint is the posterior probability $\Pr(\Delta G_{2100}>0\mid D,\mathcal{M})$, conditional on historical data $D$ and model family $\mathcal{M}$.

The registered assumptions are: common input vintage per run; pre-declared country, age, and sex harmonization; unchanged accounting logic across variants; frozen rolling holdouts; a demographic-only baseline; and aggregation from countries to regions and world. Climate-disruption features, if used, are predictive-only, lagged, and approved before fitting. No assumption converts a forecast association into a causal effect.

\subsection{Falsifiable hypotheses}

\textbf{H1:} A probabilistic component ensemble will meet pre-registered country, regional, and global calibration criteria more reliably than a deterministic central-input baseline. \textbf{H2:} Conditional component-swap diagnostics will identify stable fertility, survival, or migration contributions to peak timing across holdouts. \textbf{H3:} A climate-feature model will be retained only if it improves primary held-out scores and calibration beyond the matched demographic baseline. Each hypothesis can fail. A poorer score, miscalibrated interval, unstable attribution, or null climate ablation is an informative outcome, not a reason to modify the result after inspection.

\begin{table}[t]
\caption{Operational variables and controls}
\label{tab:variables}
\centering
\footnotesize
\begin{tabular}{L{0.23\columnwidth}L{0.61\columnwidth}}
\toprule
Element & Operational definition \\
\midrule
Independent inputs & Fertility, survival, migration trajectory family, and climate-feature set. \\
Primary endpoint & Posterior probability that $\Delta G_{2100}>0$. \\
Secondary endpoints & First-peak distribution, $G_{2100}$ distribution, interval coverage, and predictive scores. \\
Controls & Input vintage, harmonization, forecast origins, posterior-draw count, baseline priors, and aggregation map. \\
Moderators & Starting age structure, data completeness, development stratum, and approved feature lags. \\
\bottomrule
\end{tabular}
\end{table}

\section{Study Design and Methods}

\subsection{Data and preprocessing plan}

The core execution dataset will be an archived WPP 2024 release, subjected to license and version review before use. The unit of analysis is country-age-sex-year. A data manifest will record source URL, retrieval date, file checksum, release date, geographic mapping, age bins, unit conversion, exclusions, and imputation status. This is necessary because a forecast result can change when its geography or revision changes.

Supplemental climate-disruption data are not assumed available or valid by default. Before use, reviewers must approve provenance, temporal frequency, spatial mapping, country coverage, lag construction, missing-data treatment, and intended predictive-only interpretation. Data missingness will be reported by country, year, and variable. The primary analysis will not silently fill gaps; any model-based imputation is a separately registered sensitivity analysis compared with complete-case calibration subsets.

\subsection{Model comparison}

The deterministic baseline propagates a common central path through Eqs.~\eqref{eq:cohort}--\eqref{eq:global}. The demographic probabilistic ensemble generates coherent draws for fertility, survival, and migration. It may use partial pooling to regularize sparse country histories while preserving country-specific residual uncertainty. The third model is identical except for the addition of the approved climate feature set inside one or more component models. Its purpose is to test predictive increment, not causation.

All versions use the same historical cutoff, holdout window, posterior-draw budget, accounting logic, and aggregation hierarchy. This eliminates a common but avoidable comparison error: giving the more complex model a different data vintage or a longer tuning opportunity. No model is allowed to choose its own endpoint after seeing future outcomes.

\subsection{Rolling-origin validation}

At a registered origin year, observations after that origin are held out. Each model is fitted only to earlier information, projects the held-out period, and is scored against the withheld demographic outcomes. The process is repeated over at least three non-overlapping origins when data coverage permits. Chain convergence, effective sample size, and Monte Carlo error must be acceptable before an origin enters the comparison.

For a holdout $h$ with outcome $y_h$ and predictive density $p_{\mathcal{M}}$, the log predictive density is
\begin{equation}
\lpd(\mathcal{M})=\sum_{h\in\mathcal{H}}\log p_{\mathcal{M}}(y_h\mid D_{<h}),
\label{eq:lpd}
\end{equation}
where $\mathcal{H}$ is the frozen set of holdout observations and $D_{<h}$ is information available before the associated forecast origin. The score is assessed alongside continuous ranked probability score, bias, interval width, and calibration; it is not the sole decision criterion.

For nominal interval level $q$, empirical coverage is
\begin{equation}
\widehat{C}_{q}=\frac{1}{|\mathcal{H}|}\sum_{h\in\mathcal{H}}\mathbb{I}\{y_h\in I_{h,q}\},
\label{eq:coverage}
\end{equation}
where $I_{h,q}$ is the $q$-level predictive interval and $\mathbb{I}$ is the indicator function. Coverage is examined by country, age, sex, and aggregation level because a globally acceptable average can hide repeated local failure.

\subsection{Ablations, attribution, and decision}

The primary ablation compares demographic components only with the matched climate-feature model. A climate feature set is retained only after meeting a pre-registered practical improvement margin in held-out score and calibration, remaining stable over origins, and avoiding dependence on a small set of high-leverage observations. If it fails, the climate-free model remains primary and the null increment is reported.

Component-swap diagnostics exchange one aligned fertility, survival, or migration draw between posterior trajectories that have different peak-persistence outcomes. The resulting movement in $\tau^{(r)}$ and $\Delta G_{2100}^{(r)}$ is a conditional sensitivity measure. It is not a policy-intervention effect because the components are correlated in real populations.

Let $p^*$ be a threshold registered before the final horizon is examined, $K$ the calibration gates, and $S$ the required sensitivity checks. Define the posterior persistence probability as
\begin{equation}
q=\Pr(\Delta G_{2100}>0\mid D,\mathcal{M}),
\label{eq:persistence_probability}
\end{equation}
where $q$ is conditional on observed historical data $D$ and model family $\mathcal{M}$. The decision is
\begin{equation}
d=\begin{cases}
\text{persistence branch}, & q\geq p^*\ \text{and}\ K=S=1,\\
\text{earlier-peak branch}, & q<p^*\ \text{and}\ K=S=1,\\
\text{inconclusive}, & \text{otherwise}.
\end{cases}
\label{eq:decision}
\end{equation}
Equation \eqref{eq:decision} constrains interpretation. It does not determine which branch will occur.

\begin{table*}[t]
\caption{Execution stages, controls, and stopping conditions}
\label{tab:protocol}
\centering
\scriptsize
\begin{tabular}{L{0.12\textwidth}L{0.29\textwidth}L{0.19\textwidth}L{0.15\textwidth}}
\toprule
Stage & Planned action & Controls and diagnostics & Stop or fail condition \\
\midrule
1. Register inputs & Archive WPP release and candidate feature metadata. & Checksums, licenses, geography and age mapping. & Unresolved license, incompatible geography, or undocumented revisions. \\
2. Freeze validation & Set rolling origins, horizons, priors, and thresholds. & Versioned configuration and seed ledger. & Cutoff or threshold changed after holdout inspection. \\
3. Fit baseline & Fit deterministic and probabilistic demographic models. & Convergence, effective sample size, accounting coherence. & Nonconvergence or invalid cohort accounting. \\
4. Score holdouts & Calculate predictive scores and interval calibration. & Country, age, sex, region, and global diagnostics. & Calibration gate fails or baseline comparison is not reproducible. \\
5. Run ablations & Compare climate and component-swap variants. & Same input vintage, draws, and origins. & Claimed gain depends on leakage, a few cases, or unchecked missingness. \\
6. Report branch & Apply Eq.~\eqref{eq:decision} and release provenance. & Sensitivity and communication review. & Any unresolved governance or ethical-review issue. \\
\bottomrule
\end{tabular}
\end{table*}

\section{Uncertainty Decomposition and Aggregation Diagnostics}

The central scientific risk is not merely wide posterior intervals. It is the possibility that an apparent world-level persistence result is driven by a structural choice that remains invisible in an aggregate fan chart. The execution plan will therefore keep three uncertainty layers separate: stochastic or posterior uncertainty conditional on a model, uncertainty caused by component assumptions, and between-model structural uncertainty. These layers answer different questions. A narrow conditional interval cannot compensate for an untested fertility model, and a broad between-model range cannot be interpreted as a well-calibrated probability without documented model weights.

Each posterior draw will retain a provenance record containing the input release, historical cutoff, component family, prior set, feature set, random seed, and country-harmonization map. The reporting system will group draws by these labels before producing world totals. This makes it possible to display whether a persistence branch arises across several defensible families or only within one convenient specification. A result that survives only after excluding a reasonable input vintage or country mapping is reported as structurally fragile.

\subsection{Aggregation coherence and leakage tests}

The accounting identity requires that country trajectories sum to regional and world trajectories. The protocol will calculate a numerical aggregation residual at every draw and year. A non-negligible residual signals a data transformation, geography mapping, or computational error and blocks interpretation of the global endpoint. In addition, all validation transformations must be fit using data available before their corresponding forecast origin. Feature normalization, missing-data procedures, and partial-pooling hyperparameters are recorded as potential leakage points and audited with the same temporal split as the demographic forecasts.

The study will report calibration by country, age, sex, region, and world total. This granularity guards against a misleading result in which high-population countries make global coverage look acceptable while smaller or data-poor countries repeatedly fail. It also makes migration assumptions inspectable: an aggregate global migration sum may be close to zero while country-age allocation errors still materially alter fertility exposure and the timing of peaks.

\subsection{Conditional component attribution}

For a posterior pair consisting of an earlier-peak draw $E$ and a persistence draw $P$, let $z_k$ denote the trajectory of component $k$ in the set of fertility, survival, and migration. A planned swap attribution for peak timing is
\begin{equation}
A_k=\tau(z_k^P,z_{-k}^E)-\tau(z_k^E,z_{-k}^E),
\label{eq:attribution}
\end{equation}
where $z_{-k}^E$ denotes all non-$k$ components taken from draw $E$. The quantity $A_k$ is a conditional sensitivity, not a causal intervention effect: correlated demographic processes and unmodeled social mechanisms prevent that interpretation. Its role is to reveal whether the route to persistence is consistently associated with one component family or changes when the comparison pair, aggregation level, or historical cutoff changes.

\begin{table*}[t]
\caption{Pre-specified diagnostics for a credible persistence decision}
\label{tab:diagnostics}
\centering
\footnotesize
\begin{tabular}{L{0.19\textwidth}L{0.27\textwidth}L{0.45\textwidth}}
\toprule
Diagnostic & Planned calculation & Interpretation and stopping rule \\
\midrule
Temporal calibration & Rolling-origin score, bias, interval coverage, and interval width. & Do not interpret a 2100 persistence probability if key component forecasts fail their registered historical calibration gate. \\
Aggregation coherence & Country sum minus independently reported region and world totals for each draw-year. & Any material residual blocks global endpoint reporting until geography and transformation errors are resolved. \\
Structural sensitivity & Change in persistence probability and peak distribution across priors, input vintages, and approved model families. & A decision that reverses under reasonable choices is classified as inconclusive. \\
Climate increment & Difference in validated performance between the demographic baseline and matched climate-feature model. & A climate claim is not retained when increment is absent, unstable, or concentrated in leaked or high-leverage observations. \\
Attribution stability & Distribution of $A_k$ across posterior pairs and forecast origins. & Component narratives remain conditional when attribution changes sign or magnitude across registered comparisons. \\
\bottomrule
\end{tabular}
\end{table*}

\section{Expected Outcome Branches and Conditional Conclusions}

The proposal pre-registers three conclusion branches. The first is a persistence branch: a model may support the statement that its specified conditions assign substantial posterior mass to positive global growth at 2100 only if the model passes calibration and sensitivity gates. Even this branch would not establish indefinite growth; it would state a result limited to the horizon, data vintage, model family, and registered conditions.

The second is an earlier-peak branch. If a calibrated model assigns low posterior probability to positive growth at 2100 and this is stable across input, prior, and aggregation checks, the report may support an earlier peak under the registered conditions. This would not show that a peak has already been observed or that all alternative social and demographic trajectories are impossible.

The third branch is inconclusive. If the ensemble fails held-out calibration, if the complex model has no advantage over the baseline, if the climate ablation is null or unstable, or if decision status changes materially under reasonable sensitivity analysis, no post-2100 claim is warranted. An inconclusive result is scientifically preferable to a sharp but uncalibrated end-of-century headline.

Several outcomes are explicitly not expected results because this paper does not execute the protocol: no global peak year, no 2100 population number, no estimated climate coefficient, no confidence interval, no country ranking, and no policy benefit is reported. The absence of such numbers is a boundary of the proposal, not a missing calculation.

\section{Risks, Limitations, and Human Review Requirements}

\subsection{Structural and statistical risks}

Historical backtesting cannot fully validate an unprecedented future regime. Fertility convergence may change, survival trends may face innovations or shocks, migration policy can shift abruptly, and climate-disruption exposures can interact with unmeasured institutions. A well-calibrated short or medium holdout does not erase structural uncertainty at 2100. The report must therefore display posterior uncertainty, between-model dispersion, and sensitivity to prior and input choices separately rather than collapsing them into one band.

National estimates are also not neutral raw facts: geography, territories, registration completeness, and revisions are substantive modeling choices. The execution team must publish mappings and exclusions. Aggregation requires care because an apparent global result can conceal divergent national outcomes and because global net migration has different accounting meaning from country-level flows.

\subsection{Causal and ethical boundaries}

Climate features are vulnerable to confounding with development, conflict, health systems, data coverage, and adaptation. The proposed predictive ablation cannot settle causal questions. Causal language requires a different identification strategy, a separately reviewed evidence base, and explicit treatment of interventions and interference.

Population research also has a history of misuse. No result from this design may be converted into coercive reproductive policy, stigmatization of migrants, or a ranking of populations. The study contains no individual records and no human-subject intervention, but aggregate results can still cause harm through communication. Every public presentation must foreground conditionality, uncertainty, and the distinction between describing demographic possibilities and prescribing human choices.

\subsection{Required human review}

Before execution, an independent demographer must approve component definitions and country harmonization. A statistician must approve priors, convergence criteria, scoring, and the decision threshold. A climate-risk scientist must approve the feature provenance and predictive-only language. A data-governance reviewer must confirm licenses and absence of personal records. An ethics reviewer must check the final communication for coercive or stigmatizing implications. These reviews are not ceremonial: any unresolved item stops execution and limits the result to this design document.

\section{References}

\begin{thebibliography}{00}
\bibitem{unwpp2024} United Nations Department of Economic and Social Affairs, Population Division, "World Population Prospects 2024: Dataset," 2024. [Online]. Available: \url{https://www.un.org/development/desa/pd/content/world-population-prospects-2024}
\bibitem{vollset2020} S. E. Vollset \textit{et al.}, "Fertility, mortality, migration, and population scenarios for 195 countries and territories from 2017 to 2100: a forecasting analysis for the Global Burden of Disease Study," \textit{The Lancet}, vol. 396, no. 10258, pp. 1285--1306, 2020, doi: 10.1016/S0140-6736(20)30677-2.
\bibitem{gerland2014} P. Gerland \textit{et al.}, "World population stabilization unlikely this century," \textit{Science}, vol. 346, no. 6206, pp. 234--237, 2014, doi: 10.1126/science.1257469.
\bibitem{raftery2014} A. E. Raftery, L. Alkema, and P. Gerland, "Bayesian population projections for the United Nations," \textit{Statistical Science}, vol. 29, no. 1, pp. 58--68, 2014, doi: 10.1214/13-STS419.
\bibitem{azose2016} J. J. Azose, H. Sevcikova, and A. E. Raftery, "Probabilistic projections of population with migration uncertainty," \textit{Proceedings of the National Academy of Sciences}, vol. 113, no. 23, pp. 6460--6465, 2016, doi: 10.1073/pnas.1606119113.
\bibitem{raftery2023} A. E. Raftery and H. Sevcikova, "Probabilistic population forecasting: Short to very long-term," \textit{International Journal of Forecasting}, vol. 39, no. 1, pp. 73--97, 2023, doi: 10.1016/j.ijforecast.2021.09.001.
\end{thebibliography}

\appendices

\section{Idea-Evolution Record}

No true MCTS, evolutionary search, or autonomous route trajectory exists for this project. The audit record is a frozen comparison of three routes: projection reconciliation, the selected component-resolved atlas, and a rejected climate-only regression. The selected route targets accepted gaps GAP-PP-01 and GAP-PP-02. The climate-only route was rejected because the Survey registry does not verify a homogeneous causal climate-to-population mechanism. The integrated assessment route remains high risk because it would require sources and assumptions beyond the registered evidence set.

\section{Variables, Symbols, and Operational Definitions}

\begin{table}[t]
\caption{Symbols used in the proposed analysis}
\label{tab:symbols}
\centering
\footnotesize
\begin{tabular}{L{0.24\columnwidth}L{0.66\columnwidth}}
\toprule
Symbol & Definition \\
\midrule
$c,a,s,t$ & Country, age category, sex category, and year index. \\
$P_{c,a,s,t}$ & Population state in the harmonized country-age-sex cell. \\
$S_{c,a,s,t}$ & Survival probability for a cohort transition. \\
$M_{c,a,s,t}$ & Net migration assigned to the cell. \\
$f_{c,a,s,t}$ & Age-specific fertility rate. \\
$B_{c,t}$ & Births generated in country $c$ at year $t$. \\
$G_t,\Delta G_t$ & World total population and annual global change. \\
$\tau^{(r)}$ & First peak year for posterior draw $r$, with right censoring at 2100. \\
$p^*,K,S$ & Registered persistence threshold, calibration gates, and sensitivity gates. \\
\bottomrule
\end{tabular}
\end{table}

\section{Evidence Coverage and Review Checklist}

The Author stage used only citation keys registered in the Survey manifest: unwpp2024, vollset2020, gerland2014, raftery2014, azose2016, and raftery2023. The report introduces no new source, empirical result, mechanism, or citation beyond that frozen set. Claim CLM-01 supports the WPP data description; CLM-02 supports the conditional Vollset reference scenario; CLM-03 supports probabilistic component and migration methodology; CLM-04 supports the comparison of changing projections; and CLM-05 limits climate language to a predictive design inference.

The completed execution package must additionally verify data-release licenses, register exact holdout cutoffs and thresholds, archive code and seeds, audit geographic harmonization, inspect convergence and calibration, and complete the listed demographic, statistical, climate-risk, data-governance, and ethics reviews. Until then, this document remains a proposal with no observed results.

\end{document}
